Reservoir computing was introduced to exploit high-dimensional recurrent dynamics while training only a simple readout~\cite{maass2002real,jaeger2004harnessing,lukosevicius2009reservoir}. ESNs remain the most direct classical comparator for QRC because they share the same learning structure: fixed recurrent dynamics, nonlinear temporal features, and a trained linear readout. NVAR provides an important fixed-feature robustness baseline for nonlinear forecasting, but it is not the closest reservoir-to-reservoir comparison~\cite{gauthier2021next}.

QRC transfers this architecture to quantum dynamics. Disordered spin ensembles, spatial multiplexing, and driven quantum systems have been proposed as fixed temporal feature maps~\cite{fujii2017harnessing,nakajima2019boosting,ghosh2019quantum}. More recent work shows why protocol details matter: spin-network performance depends on dynamical regimes~\cite{martinezpena2021dynamical}, dissipation can be an active resource rather than a nuisance~\cite{sannia2024dissipation}, and input encoding can itself provide strong nonlinear transformations~\cite{govia2022nonlinear}. These results motivate our split between mechanistic, encoding-enhanced, and dissipative QRC variants.

Time-series benchmark archives and forecasting competitions show that model ranking is strongly dataset-dependent: seasonality, persistence, noise, nonstationarity, and regime changes can dominate average scores~\cite{makridakis2018m4,godahewa2021monash}. Our contribution is to connect this dataset-regime view to QRC. Rather than reporting only a mean QRC-versus-ESN score, we measure dataset properties, label empirical advantage under a frozen protocol, and learn where in property space QRC becomes useful.

The paper is therefore complementary to mechanism studies. It identifies \emph{where} QRC is empirically worth testing under a controlled benchmark. It does not prove \emph{why} those regimes arise from a particular quantum mechanism.
